\documentclass[12pt,a4paper]{article}

\usepackage[utf8]{inputenc}
\usepackage[T1]{fontenc}
\usepackage[turkish]{babel}
\usepackage{geometry}
\usepackage{hyperref}
\usepackage{listings}
\usepackage{xcolor}
\usepackage{enumitem}
\usepackage{booktabs}
\usepackage{titlesec}
\usepackage{fancyhdr}
\usepackage{graphicx}

\geometry{margin=2.5cm}

\definecolor{codegreen}{rgb}{0,0.6,0}
\definecolor{codegray}{rgb}{0.5,0.5,0.5}
\definecolor{codepurple}{rgb}{0.58,0,0.82}
\definecolor{backcolour}{rgb}{0.97,0.97,0.97}
\definecolor{primaryviolet}{RGB}{139,92,246}

\lstdefinestyle{tsstyle}{
    backgroundcolor=\color{backcolour},
    commentstyle=\color{codegreen},
    keywordstyle=\color{magenta},
    numberstyle=\tiny\color{codegray},
    stringstyle=\color{codepurple},
    basicstyle=\ttfamily\footnotesize,
    breakatwhitespace=false,
    breaklines=true,
    captionpos=b,
    keepspaces=true,
    numbers=left,
    numbersep=5pt,
    showspaces=false,
    showstringspaces=false,
    showtabs=false,
    tabsize=2,
    language=python
}
\lstset{style=tsstyle}

\pagestyle{fancy}
\fancyhf{}
\lhead{CodeAlchemist Staj Raporu}
\rhead{Hafta 10}
\cfoot{\thepage}

\title{\textbf{CodeAlchemist Geliştirme Raporu\\Hafta 10: Güvenlik, Ödeme Sistemleri ve Teknik Detaylar}}
\author{Dilakemer}
\date{Nisan - Mayıs 2026}

\begin{document}

\maketitle

\section{Hafta Özeti}

Bu hafta CodeAlchemist platformunda güvenlik, yerel ödeme sistemleri ve VS Code eklentisi kullanıcı deneyimi üzerine yoğunlaşıldı. Yapılan geliştirmeler, sistemin hem ticari hem de teknik güvenliğini bir üst seviyeye taşımıştır.

\begin{enumerate}
    \item \textbf{Iyzico Entegrasyonu} -- Yerel pazar desteği için TRY para birimi ve Iyzico API altyapısı kuruldu.
    \item \textbf{Command Card Mekanizması} -- Eklenti tarafında komut onay süreçleri görselleştirildi.
    \item \textbf{API Key Hashing} -- Veritabanı sızıntılarına karşı API anahtarları SHA-256 ile koruma altına alındı.
    \item \textbf{Token Senkronizasyonu} -- Web ve Extension arasında kesintisiz bakiye aktarımı ve harcama kontrolü sağlandı.
\end{enumerate}

\section{Teknik Detaylar ve Uygulama}

\subsection{API Key Hashing ve Güvenlik}

Sistemde kullanılan API anahtarları artık düz metin olarak saklanmamaktadır. Her anahtar oluşturulduğunda veya migrasyon sırasında hashlenir.

\begin{lstlisting}[caption={API Key Hashing Mantığı (app.py)}]
def _hash_api_key_value(raw_key: str) -> str:
    return hashlib.sha256((raw_key or '').encode('utf-8')).hexdigest()

def _build_stored_api_key(raw_key: str) -> str:
    token = (raw_key or '').strip()
    client = _detect_api_key_client(token)
    prefix = 'ca-vsc' if client == 'vscode' else 'ca-web'
    return f"{prefix}-h-{_hash_api_key_value(token)}"
\end{lstlisting}

Ayrıca, veritabanı başlatıldığında mevcut anahtarlar için otomatik bir migrasyon scripti çalıştırılmaktadır.

\subsection{Iyzico Ödeme Sistemi}

Iyzico entegrasyonu, Türkiye'deki kullanıcıların TRY üzerinden token satın alabilmesine olanak tanır. Backend tarafında Iyzico Checkout Form Initialize Request yapısı kullanılmıştır.

\begin{lstlisting}[caption={Iyzico Checkout Session Oluşturma (app.py)}]
@app.route('/api/billing/iyzico/checkout-session', methods=['POST'])
def iyzico_checkout_session():
    # ... paket ve kullanıcı doğrulama ...
    request_data = {
        'locale': 'tr',
        'price': str(pkg.price_try),
        'paidPrice': str(pkg.price_try),
        'currency': 'TRY',
        'basketId': 'B' + str(int(time.time())),
        'paymentGroup': 'PRODUCT',
        'callbackUrl': request.host_url.rstrip('/') + '/api/billing/iyzico/callback',
        # ... kullanıcı bilgileri ...
    }
    checkout_result = iyzipay.CheckoutFormInitialize().create(request_data, iyzico_options)
    return jsonify({'payment_url': checkout_result.get('paymentPageUrl')})
\end{lstlisting}

\subsection{Command Card ve Extension UX}

Agent tarafından önerilen terminal komutları, kullanıcıya interaktif bir kart yapısında sunulur.

\begin{lstlisting}[language=html, caption={Command Card Rendering (chatView.ts)}]
function addCommandCard(action) {
    const card = document.createElement('div');
    card.className = 'command-card';
    card.innerHTML = `
        <div class="command-header"><span>Terminal Komutu</span></div>
        <div class="command-body">
            <div class="command-text">${escapeHtml(action.command)}</div>
        </div>
        <div class="command-actions">
            <button class="btn btn-primary run-btn">Onayla</button>
            <button class="btn btn-secondary discard-btn">Reddet</button>
        </div>
    `;
    // ... event listeners ...
}
\end{lstlisting}

\subsection{Bakiye Senkronizasyonu ve OTP Geçişi}

Eklenti içerisinden bakiye yüklemek istendiğinde, web tarayıcısına güvenli bir geçiş yapmak için OTP (One-Time Password) mekanizması kullanılır.

\begin{lstlisting}[caption={OTP ile Güvenli Yönlendirme (chatProvider.ts)}]
private async _handleOpenPurchase() {
    const otpRes = await getVscodeOtp(endpoint, apiKey);
    if (otpRes.status === 'ok') {
        const secureUrl = `${this._purchaseUrl}?auth_otp=${otpRes.otp}`;
        await vscode.env.openExternal(vscode.Uri.parse(secureUrl));
    }
}
\end{lstlisting}

\section{Etkilenen Dosyalar}

\begin{center}
\begin{tabular}{@{}lll@{}}
\toprule
\textbf{Dosya} & \textbf{Değişiklik Türü} & \textbf{Açıklama} \\
\midrule
\texttt{server/app.py}        & Güncellendi & Iyzico rotaları, Hashing mantığı, OTP üretimi \\
\texttt{server/models.py}     & Güncellendi & \texttt{price\_try}, \texttt{VSCodeOTP} modelleri \\
\texttt{src/chatProvider.ts}  & Güncellendi & Token Guard, OTP transfer mantığı \\
\texttt{src/chatView.ts}      & Güncellendi & Command Card, Payment Card UI ve CSS \\
\bottomrule
\end{tabular}
\end{center}

\section{Doğrulama}

\begin{enumerate}
    \item \textbf{Güvenlik:} Veritabanındaki tüm anahtarların hashlenmiş olduğu manuel SQL sorguları ile doğrulandı.
    \item \textbf{Ödeme:} Iyzico sandbox ortamında ödeme tamamlandığında bakiyenin extension tarafında anlık (WebSocket/Polling) güncellendiği görüldü.
    \item \textbf{UX:} Command Card üzerinden komut onaylandığında, VS Code terminalinin doğru CWD ile açıldığı test edildi.
\end{enumerate}

\section{Sonuç}

Bu hafta tamamlanan özellikler ile CodeAlchemist platformu ticari bir ürün olma yolunda önemli bir eşiği aşmıştır. Özellikle güvenlik ve yerel ödeme kanalları, kullanıcı güvenini artıran kritik unsurlardır.

\vspace{1cm}
{\color{primaryviolet}\rule{\linewidth}{1pt}}
\begin{center}
\textit{CodeAlchemist Staj Raporu -- Hafta 10 -- Mayıs 2026}
\end{center}

\end{document}